\documentclass[11pt]{article}
\usepackage[T1]{fontenc}
\usepackage{textcomp}
\usepackage[margin=0.75in]{geometry}
\usepackage{amsmath,amssymb,amsthm}
\usepackage{array,booktabs}
\usepackage{hyperref}
\setlength{\parindent}{0pt}
\newtheorem{theorem}{Theorem}
\theoremstyle{definition}
\newtheorem{definition}{Definition}
\title{Flow Proof Artifact: prop-07-pythagoras-on-diameter}
\author{Generated by \texttt{flow doc proof}}
\date{Flow Proof Book}
\begin{document}
\maketitle
If on the diameter of a circle a point is taken and a perpendicular is erected, the square on the half-diameter exceeds the square on the line to the point by the square on the perpendicular.\\[0.5em]
\textbf{Source.} euclid — Elements, Book III, Proposition 7\\[0.5em]
\bigskip
\noindent{\large \textbf{Derived fact 1.} \textit{If on the diameter of a circle a point is taken and a perpendicular is erected, the square on the half-diameter exceeds the square on the line to the point by the square on the perpendicular} \hfill the Euclidean plane \cdot Euclid Book III \cdot Proposition 7: the square on a half-diameter exceeds the square on a line to an interior point}\par
\smallskip\noindent\textit{Coordinate: the Euclidean plane \cdot Euclid Book III \cdot Proposition 7: the square on a half-diameter exceeds the square on a line to an interior point}\par
\smallskip\noindent\textit{Source: euclid — Elements, Book III, Proposition 7}\par
\smallskip\noindent\textit{Built on: proposition 47: square on hypotenuse equals sum of squares on legs, for Book I of the Elements in on the Euclidean plane}\par
\medskip
\noindent\textbf{Goal.} If on the diameter of a circle a point is taken and a perpendicular is erected, the square on the half-diameter exceeds the square on the line to the point by the square on the perpendicular\par
\noindent$\displaystyle sq(\text{half diameter}) = sq(\text{line to point}) + sq(perpendicular)$\par
\medskip
\noindent
\begin{tabular}{@{} >{\raggedright\arraybackslash}p{0.06\textwidth} >{\raggedright\arraybackslash}p{0.47\textwidth} >{\raggedleft\arraybackslash}p{0.41\textwidth} @{}}
\textbf{\#} & \textbf{Proof} & \textbf{Mathematics} \\
\midrule
\textbf{1.} & We prove that proposition 7: the square on a half-diameter exceeds the square on a line to an interior point for Euclid Book III the Euclidean plane. & \\[0.45em]
\textbf{2.} & Let AB = diameter. & \\[0.45em]
\textbf{3.} & Let C = interior\_point. & \\[0.45em]
\textbf{4.} & Let D = foot\_of\_perpendicular. & \\[0.45em]
\textbf{5.} & From step 2, step 3, and step 4, we can deduce that sq(half) equals sq(AC) plus sq(CD). & $\displaystyle sq(half) = sq(AC) + sq(CD)$ \\[0.45em]
\textbf{6.} & We invoke the derived fact governing Book I of the Elements in the Euclidean plane: proposition 47: square on hypotenuse equals sum of squares on legs, for Book I of the Elements in on the Euclidean plane. & \\[0.45em]
\textbf{7.} & From step 6, this implies sq(half diameter) equals sq(line to point) plus sq(perpendicular). Hence proven. & $\displaystyle sq(\text{half diameter}) = sq(\text{line to point}) + sq(perpendicular)$ \\[0.45em]
\bottomrule
\end{tabular}
\label{thm:Geometry-Euclid Book III-proposition_7_the_square_on_a_half_diameter_exceeds_the_square_on_a_line_to_an_interior_point}
\par\medskip\hrule
\end{document}
